\section{Métricas de Calidad}
\label{sec:metricas}

La guía de trabajo práctico exige proponer un set de métricas cuantitativas pertinentes para la norma ISO/IEC asignada, detallando sus respectivas ecuaciones matemáticas de obtención \cite{iso33000}. A continuación se presentan las métricas definidas para evaluar el cumplimiento de los controles A.14 de ISO/IEC 27001.

\subsection{Métrica 1: Densidad de Vulnerabilidades}

Mide la cantidad de vulnerabilidades encontradas por unidad de código. Es un indicador directo de la calidad de seguridad de la implementación.

\begin{equation}
D_v = \frac{V_{encontradas}}{KLOC} \times 1000
\end{equation}

Donde:
\begin{itemize}
    \item $D_v$ = Densidad de vulnerabilidades (vuln/KLOC).
    \item $V_{encontradas}$ = Número total de vulnerabilidades encontradas por SAST + DAST.
    \item $KLOC$ = Miles de líneas de código fuente.
\end{itemize}

\textbf{Interpretación:}
\begin{table}[h]
\centering
\begin{tabular}{|c|l|}
\hline
\textbf{Rango $D_v$} & \textbf{Nivel de riesgo} \\
\hline
$D_v < 1$ & Bajo \\
$1 \leq D_v < 5$ & Medio \\
$5 \leq D_v < 10$ & Alto \\
$D_v \geq 10$ & Crítico \\
\hline
\end{tabular}
\caption{Umbrales de densidad de vulnerabilidades}
\label{tab:umbral-vuln}
\end{table}

\textbf{Ejemplo de cálculo:} Si un proyecto tiene 10 KLOC y se encuentran 4 vulnerabilidades:
\begin{equation}
D_v = \frac{4}{10} \times 1000 = 400 \text{ vuln/KLOC} \quad \Rightarrow \quad \text{Nivel: Bajo}
\end{equation}

\subsection{Métrica 2: Cobertura de Análisis SAST}

Mide el porcentaje de archivos del proyecto que fueron analizados por herramientas SAST.

\begin{equation}
C_{SAST} = \frac{A_{escaneados}}{A_{total}} \times 100
\end{equation}

Donde:
\begin{itemize}
    \item $C_{SAST}$ = Cobertura de análisis SAST (\%).
    \item $A_{escaneados}$ = Archivos de código fuente analizados.
    \item $A_{total}$ = Total de archivos de código fuente.
\end{itemize}

\textbf{Objetivo mínimo:} $C_{SAST} \geq 80\%$

\textbf{Ejemplo de cálculo:} Si hay 60 archivos de código fuente y se escanean 52:
\begin{equation}
C_{SAST} = \frac{52}{60} \times 100 = 86.7\% \quad \Rightarrow \quad \text{Cumple el objetivo}
\end{equation}

\subsection{Métrica 3: Tiempo Medio de Remediación}

Mide el tiempo promedio que tarda el equipo en corregir una vulnerabilidad encontrada.

\begin{equation}
T_{rem} = \frac{\sum_{i=1}^{n} (t_{corrección_i} - t_{detección_i})}{n}
\end{equation}

Donde:
\begin{itemize}
    \item $T_{rem}$ = Tiempo medio de remediación (días).
    \item $t_{corrección_i}$ = Fecha de corrección de la vulnerabilidad $i$.
    \item $t_{detección_i}$ = Fecha de detección de la vulnerabilidad $i$.
    \item $n$ = Total de vulnerabilidades corregidas.
\end{itemize}

\textbf{Umbrales:}
\begin{table}[h]
\centering
\begin{tabular}{|c|l|}
\hline
\textbf{Tiempo} & \textbf{Categoría} \\
\hline
$T_{rem} \leq 1$ día & Crítica (corrección inmediata) \\
$1 < T_{rem} \leq 7$ días & Alta prioridad \\
$7 < T_{rem} \leq 30$ días & Prioridad media \\
$T_{rem} > 30$ días & Requiere gestión \\
\hline
\end{tabular}
\caption{Umbrales de tiempo de remediación}
\label{tab:umbral-rem}
\end{table}

\subsection{Métrica 4: Tasa de Aceptación de Seguridad}

Mide el porcentaje de pruebas de seguridad que cumplen con los criterios de aceptación.

\begin{equation}
R_{acep} = \frac{P_{aprobadas}}{P_{total}} \times 100
\end{equation}

Donde:
\begin{itemize}
    \item $R_{acep}$ = Tasa de aceptación de seguridad (\%).
    \item $P_{aprobadas}$ = Pruebas de seguridad que cumplieron criterios.
    \item $P_{total}$ = Total de pruebas de seguridad ejecutadas.
\end{itemize}

\textbf{Objetivo mínimo:} $R_{acep} = 100\%$ (todas las pruebas críticas deben aprobar antes de producción).

\subsection{Métrica 5: Nivel de Madurez del SSDLC}

Mide el nivel de madurez del proceso de seguridad en el ciclo de vida del desarrollo, basado en la escala de ISO/IEC 33000 \cite{iso33000}:

\begin{table}[h]
\centering
\small
\begin{tabular}{|c|l|p{6cm}|}
\hline
\textbf{Nivel} & \textbf{Nombre} & \textbf{Criterio general} \\
\hline
0 & Incompleto & No hay proceso de seguridad definido. \\
1 & Realizado & Se realizan pruebas de seguridad de forma ad-hoc. \\
2 & Gestionado & Hay un proceso documentado y se ejecuta de forma controlada. \\
3 & Establecido & El proceso está definido, documentado y se aplica orgánicamente. \\
4 & Predecible & Las métricas de seguridad se miden y controlan cuantitativamente. \\
5 & Optimizado & El proceso se mejora continuamente según métricas y retroalimentación. \\
\hline
\end{tabular}
\caption{Escala de madurez SSDLC según ISO/IEC 33000}
\label{tab:madurez}
\end{table}

\textbf{Ecuación de evaluación de madurez:}
\begin{equation}
M_{SSDLC} = \frac{\sum_{i=1}^{5} P_i \times w_i}{\sum_{i=1}^{5} w_i}
\end{equation}

Donde:
\begin{itemize}
    \item $P_i$ = Puntaje obtenido en la práctica $i$ (0 o 1).
    \item $w_i$ = Peso de la práctica $i$ (según prioridad).
\end{itemize}

\subsection{Resumen de Métricas}

\begin{table}[h]
\centering
\small
\begin{tabular}{|l|l|l|}
\hline
\textbf{Métrica} & \textbf{Fórmula} & \textbf{Objetivo} \\
\hline
Densidad de vulnerabilidades & $D_v = V / KLOC \times 1000$ & $D_v < 1$ \\
Cobertura SAST & $C = A_{esc} / A_{tot} \times 100$ & $C \geq 80\%$ \\
Tiempo de remediación & $T_{rem} = \Sigma \Delta t / n$ & $T_{rem} \leq 7$ días \\
Tasa de aceptación & $R = P_{apr} / P_{tot} \times 100$ & $R = 100\%$ \\
Madurez SSDLC & $M = \Sigma P_i w_i / \Sigma w_i$ & Nivel $\geq 2$ \\
\hline
\end{tabular}
\caption{Resumen de métricas de calidad para A.14}
\label{tab:resumen-metricas}
\end{table}
